% !TEX program = xelatex
% Chladni 板逆向设计系统：改进方向（终版）
% 吸收 2026-06 决策实验勘误（reports/future_directions_for_successors.zh-CN.md）
% 编译：xelatex improvement_directions_final_zh.tex
\documentclass[11pt,a4paper]{article}

\usepackage[margin=2.5cm]{geometry}
\usepackage{amsmath,amssymb}
\usepackage{booktabs}
\usepackage{enumitem}
\usepackage{array}
\usepackage{xeCJK}
\setCJKmainfont{Droid Sans Fallback}[AutoFakeBold=2.5,AutoFakeSlant=0.15]
\setCJKsansfont{Droid Sans Fallback}[AutoFakeBold=2.5]
\setCJKmonofont{Droid Sans Fallback}
\usepackage[colorlinks=true,linkcolor=blue,urlcolor=blue]{hyperref}

\newcommand{\code}[1]{\texttt{#1}}
\linespread{1.18}
\setlength{\parskip}{0.35em}

\title{\textbf{Chladni 板逆向设计系统：改进方向（终版）\\
\large 从"仿真内自证"走向"现实可复现"---吸收 2026-06 决策实验勘误，含试错史回顾}}
\author{Chladni Studio 项目组}
\date{2026 年 6 月}

\begin{document}
\maketitle

\begin{abstract}
我们的软件做的事情可以用一句话概括：用户画一个图案，程序反过来设计一块厚薄不均的板子和几个驱动频率，让沙子在振动的板上尽量摆出这个图案。本报告的初版提出了五个改进方向，重心放在优化器与评分一侧。随后的一轮决策实验（见 \code{reports/future\_directions\_for\_successors}）推翻了这个排序：项目真正的瓶颈不是算法不够聪明，而是\textbf{整条链路始终在仿真内部自我验证}---surrogate 向 COMSOL 看齐、COMSOL 被当成真理，却从未用一块真实的板标定过；同时优化器被物理天花板逼到了\textbf{物理上不可复现的离共振工作点}上。"COMSOL 里有图、现实里没有"不是某个 bug，而是这套方法论的必然结果。本终版据此重排改进方向：(0)~先闭合"仿真到真实板"的验证回路；(1)~把物理可复现性当硬约束，频率只从隔离良好的本征频率里选；(2)~诚实面对对称天花板：要么限制目标，要么改变物理；(3)~统一 surrogate、COMSOL 与真实台子的激励模型；(4)~评分与几何感知精修排在最后。初版的"中心伪影修正"已结案：机理诊断正确，但它只是画面问题，排除中心的处理早已在生产配置中生效，而初版提出的"绝对位移正修"经实测无效，不应再投入。文末保留并更新了试错史回顾。
\end{abstract}

\section{背景：系统怎么工作，以及哪里出了问题}

\subsection{流水线}

整条流水线分四步。第一步，用户给一张 $256\times 256$ 的黑白目标图。第二步，一个我们自己写的"草稿物理引擎"（业内叫 surrogate，替代模型）快速迭代约 300 步，同时调整板子的 $15\times 15$ 厚度分布 $H$ 和 6 个候选驱动频率。第三步，把板子交给精确但缓慢的 COMSOL 仿真，算出前 30 个共振模态，挑出对目标贡献最大的几个频率（Phase~1），外加一个靠扫频发现的"幸运频率"（165~Hz）。第四步（Phase~2），拿 COMSOL 结果反向微调板子，小步快试一轮。最终在 IC 字母目标上报告了"目标线上的沙密度是全板平均的 4.85 倍"。

\subsection{决策实验揭示的两个事实}

结案前的一轮只读决策实验（脚本 \code{scripts/\_decexp\_*.py}，输出在 \code{reports/\_decision\_experiments/}）给出了两个改变全局的事实：

\textbf{事实一：链路从未接触现实。}整个项目没有任何一步验证过 COMSOL 能预测真实的板。surrogate 向 COMSOL 校准、COMSOL 被当成真理，但这位"裁判"自己从未被检验过。因此 $4.85\times$ 是一个\textbf{仿真内部的数字}，从未锚定到现实。

\textbf{事实二：设计系统性地工作在离共振点上。}对 32 个"设计与 COMSOL 本征频率"配对的普查显示：23 个设计把超过一半的驱动能量打在共振点之外。基础生产设计只有 \textbf{4\% 的能量 on-resonance}：权重最大的 165.4~Hz（占 0.604，即"幸运频率"）偏离最近本征模 2.87 个半功率带宽；1099~Hz（占 0.289）偏离 25 个带宽；约 36\% 的能量落在 800--1000~Hz 这段\textbf{一个本征模都没有的"模态荒漠"}。离共振意味着强迫响应弱、多模混叠、图案淡；而真实板的模态因材料与打印偏差还会再移位---同一组频率在真实板上只会得到另一团弱响应。\textbf{所以"现实复现不出 COMSOL 的图"不是意外，是必然。}

\subsection{为什么优化器会奔向离共振？}

这正是物理天花板在起作用。单点中心激励加正方形对称（D4）意味着：\textbf{只有对称的振型家族（A1）能被有效激发且可复现}。IC 这种非对称目标在对称模态里根本表达不出来（它只有约 42\% 的成分是物理允许的；X 形是 100\%，单对角线约 52\%）。优化器按"图像像不像"打分，于是被结构性地推向那些"碰巧像 IC"的离共振工作点---而这些点恰恰不可复现。\textbf{换句话说，流水线"能做出 IC"，靠的就是物理上站不住的频率。}初版报告承认物理天花板，但严重低估了它的中心地位：它不是背景板，而是整个矛盾的发动机。

\subsection{证据的可信度分级}

吸收勘误文档的做法，本文所有关键结论按来源分级，请读者（和接手者）不要把估计当定论：

\begin{center}
\small
\begin{tabular}{@{}llp{0.5\textwidth}@{}}
\toprule
标记 & 可信度 & 含义 \\
\midrule
【COMSOL真值】 & 高 & 来自 COMSOL 本征频率 / 强迫响应导出 \\
【代码核实】 & 高 & 直接阅读源码确认 \\
【surrogate估计】 & 中 & 用 $25\times 25$ 粗网格替代模型算出，需复核 \\
【待实测】 & 未验证 & 必须用真实板实验确认 \\
\bottomrule
\end{tabular}
\end{center}

上文 4\%、23/32、模态荒漠为【COMSOL真值】；A1 可达性比例为【代码核实】（数学推导）+【待实测】（真实板上的复现性）。

\subsection*{名词速查}

\begin{center}
\small
\begin{tabular}{@{}p{0.26\textwidth}p{0.66\textwidth}@{}}
\toprule
术语 & 大白话解释 \\
\midrule
surrogate（替代模型） & 自写的快速近似物理引擎：算得糙但快，是指南针不是裁判 \\
enrichment（富集度） & 目标线上的沙密度是全板平均的几倍 \\
recall / precision & 目标线被"安静区"盖住多少 / 安静区有多大比例真的落在目标上 \\
节线 / 安静区 & 板上振幅接近零的地方，沙子聚集在这里 \\
D4 对称 / A1 振型 & 正方形的 8 种自我重合操作 / 在所有操作下都不变的最对称振型家族 \\
on-resonance & 驱动频率落在某个本征频率的半功率带宽之内，响应强且可复现 \\
半功率带宽 & 共振峰的"有效宽度"$\approx \zeta f$（阻尼比乘频率）；离开它响应骤降 \\
模态隔离度 & 所选本征频率与邻近模态的间距；隔离差则图案对误差敏感 \\
动态放大倍率 & 共振响应比直接推一下大多少倍；放大低则图案淡、基底项显眼 \\
模态荒漠 & 一段没有任何本征频率的频带；把能量打进去只能得到弱响应 \\
PCA 设计流形 & 从历史好设计中提炼出的几十个"有效改法"方向 \\
最优传输（OT） & 把当前沙图"搬"成目标图的最小力气；搬运方案告诉你沙该往哪挪 \\
\bottomrule
\end{tabular}
\end{center}

\section{方向 0（地基）：闭合"仿真到真实板"的验证回路}

接手后的第一件事，不是改任何算法，而是\textbf{打印一块板、实际测它}：用敲击法或激光测振测本征频率与模态（最低成本的替代：扫频、撒沙、拍照），再反过来用实测数据拟合 COMSOL 模型的材料与边界参数。

为什么它排第零位：在它完成之前，COMSOL 是一位\textbf{从未被检验过的裁判}，所有下游数字（包括 $4.85\times$）都悬在空中。其余一切改进---更好的评分、更聪明的搜索---都建立在这位裁判可信的前提上。\textbf{没有这一环，其余改进都是空中楼阁。}

具体可执行的最小版本：(a)~打印当前 production\_design 与一块均匀板；(b)~扫频记录沙图随频率的变化（手机慢动作即可起步）；(c)~把实测本征频率与 COMSOL 预测列表对照，拟合杨氏模量、密度与夹持刚度；(d)~把拟合后的参数写回 \code{config.yaml}，重新跑一遍 Phase~1，看模态预测是否对得上。这一步同时回答"COMSOL 和现实差多少"与"差距主要来自哪个参数"。【待实测】

\section{方向 1：把"物理可复现性"当硬约束（吸收原"频率内循环"）}

\subsection{约束本身}

证据见 1.2 节。新规则很简单：\textbf{只在隔离良好的本征频率上驱动}（与最近邻模态的间距远大于半功率带宽，这样即使真实板和仿真有偏差，工作点也不会掉出共振）；并在优化目标里显式惩罚离共振与密集模态区。建议引入"可复现性分"：
\begin{equation}
S_{\mathrm{repro}} \;=\; \text{动态放大倍率} \times \text{模态隔离度},
\end{equation}
低于阈值的设计，无论图案多像目标，一律淘汰。对当前 production\_design，按此标准推荐的候选频率为 \textbf{58 / 132 / 264 / 371 Hz}（$\le 600$~Hz、间距 $\ge 4$ 个带宽），并避开 800--1000~Hz 模态荒漠。【surrogate估计，需 COMSOL 复核】

\subsection{原"频率内循环"的正确归宿}

初版报告提出"每步先给板子扫全频段、在最佳频率处评分再求梯度"（用模态叠加法让扫频几乎免费、按 Danskin 定理在最优频率处求梯度）。这个想法在新框架下\textbf{不是被废弃，而是被收编}：一旦频率只允许从 COMSOL 本征频率集合里选，"连续频段上取 max"就退化为\textbf{"在本征频率列表上做离散选择 + 隔离度惩罚"}---这恰好和可复现性约束是同一件事。每步对 proxy 板做一次特征分解，取其本征频率列表（而非连续扫频），按"图案得分 $\times$ 可复现性分"选驱动点，再对厚度求梯度。原方案里"步长要小、让新板的最优频率与旧板可比"的信任域直觉原样保留。

\section{方向 2：诚实面对天花板---要么换目标，要么换物理}

物理事实（1.3 节）摆在那里，只有两条出路，必须二选一或并行：

\textbf{(a) 换目标}：把目标库限制在\textbf{物理可达的图案}上---A1 对称成分高的图形（X、十字、环、星形），或从均匀板模态库（\code{src/uniform\_pattern/}）出发做小幅微扰的"模态邻域"图案。配套动作是把已有的可实现性裁决（\code{target\_realizability.py}、D4 分解的可达性分数）\textbf{前置到 UI 入口}：用户画完图，先告诉他"这幅图有多少成分是物理允许的"，低于阈值就引导改图，而不是烧两小时算力后给一张不像的结果。

\textbf{(b) 换物理}：真正打破 D4 对称约束---多激振器相控阵、分布式或边缘激励、非对称板形。被搁置的 \code{mosaic\_z} 六点相控方案其实是这条路的正确雏形（前提：真实台子做得出多点相控驱动）。这是"想做出任意图案"唯一物理上诚实的路线。【待实测】

科普版结尾那句话在技术上完全成立："别再从中心摇这块板了。多加几个激振器，或者从偏一侧去推，或者换一种板形---一整片全新的图案世界就会打开。"

\section{方向 3：统一三处激励模型}

【代码核实】当前 surrogate 与 COMSOL 用的不是同一个物理问题：surrogate 是\textbf{整板基础惯性激励}（\code{direct\_forced\_response}：$f=-M\,a_{\mathrm{base}}$，相当于整块板坐在振动台上），而 COMSOL runner 是\textbf{带位置、幅值、相位的高斯点激励}（\code{run\_chladni\_forced\_response.m}）。优化器在物理 A 下选板子，却在物理 B 下验收---这层错配一直是隐患。真实台子用的是单个中心激振器。

要做的事：把 surrogate、COMSOL、真实台子三处的激励模型统一成同一种描述（建议以真实台子为准：中心点激励 + 实测激励幅值），并做一次"surrogate 预测 vs COMSOL 结果"的系统对照，量化 surrogate 到底在哪些频段、哪些指标上可信。这也是让方向 1 的"可复现性分"在 surrogate 阶段就能算准的前提。

\section{方向 4（最后）：评分与几何感知精修}

初版报告的方向一、二、四（评分重构、几何损失、随机重启）的内容本身\textbf{没有错}，但次序错了：\textbf{在未校准、离共振的地基上把评分做得更精，只是把一个虚构优化得更精确。}等方向 0--3 落地、裁判可信之后，这些手段才有意义。届时按原方案执行：

\begin{enumerate}[leftmargin=2em]
\item \textbf{及格线 + "像不像"排序}：富集度与对比度退为可行性门槛（hinge 罚保持可导）；排序键改用低谷掩膜与目标的 $F_\beta$ 分数（同时看 recall 与 precision，杜绝"X 形拿对角线高分"的假阳性）。顺带修掉"全零振幅场 recall = 1.0"的边界 bug。
\item \textbf{几何感知损失}：距离场（SDF）与最优传输（Sinkhorn）损失让优化器知道"沙偏在哪、往哪挪"，配多尺度模糊；代码库已有两个半成品（\code{signed\_distance\_loss.py}、\code{sinkhorn\_loss.py}）。拓扑体检（连通块数、骨架分叉数）用于最终排序。
\item \textbf{随机重启}：恢复随机厚度初始化（当前优化器是确定性的，多起点为空操作）；在 PCA 设计流形坐标中做远距采样；重启基准从可实现性分析推导、按目标自适应、设次数上限。
\end{enumerate}

\section{已结案：中心伪影（"木星"）---一个有教育意义的勘误}

\subsection{结案结论}

仿真图中心那团永恒的沙堆（团队戏称"木星"）的\textbf{机理诊断是正确的}：中心 8~mm 被夹持、自由度被置零（【代码核实】\code{orthotropic\_plate.py}），撒粉模型把"振幅小"判为"积沙"，于是只要画的是相对位移场，木星就一定在，与是否共振无关。但三个事实让它结案：

\begin{enumerate}[leftmargin=2em]
\item 现实里板中心被\textbf{螺栓和夹具占住}，本来就不会有沙图案---木星是\textbf{画面伪影}，不影响真正可用的外围图案。
\item 排除中心区的处理\textbf{早已在生产配置中生效}（\code{remove\_center\_region: true}，评分时挖掉中心），做法正确，初版报告所谓"快修"实际已在线上。
\item 初版报告提出的"正修"---改用绝对位移 $u_{\mathrm{abs}} = u_{\mathrm{rel}} + a_{\mathrm{base}}/\omega^2$、声称幻影沙会自动消失---\textbf{经实测无效}：在真实 production\_design 的前向解上，中心撒粉占比反而从 31\% 升到 46\%，外围图案被改动 14\%。原因：该设计的动态放大倍率只有约 4，基底项与响应同量级，绝对位移的"安静区"变成了 $u_{\mathrm{rel}}\approx -u_{\mathrm{base}}$ 的抵消带，与节线不重合。【surrogate估计】\textbf{不要再在这个方向上花时间。}
\end{enumerate}

\subsection{为什么保留这一节}

因为它是一个现成的方法论案例：初版报告在这里犯的，正是全项目的同款错误---\textbf{在仿真推理里给现实开药方，而没有先做实验}。"正修"在纸面上物理正确（沙确实听绝对加速度的），但在这块板的实际参数区（低放大倍率）下失效。它示范了为什么本终版把"先实测、再开方"立为第一原则。

\section{试错史回顾与勘误}

项目的试错记录（trails）保存了很多有价值的判断，这一节把它们和新证据接到一起。

\subsection{材料之路：从"均匀"到"各向异性"，再到撤销 $\theta$}

故事线：先用 SLA 打印追求均匀，发现厚度差异再大、图案仍顽固对称；改用 FDM 各向异性材料；先让每格纤维取向 $\theta$ 不同，效果不好，退回全板同一取向。

背后的物理：纯靠厚度破对称很弱（弯曲刚度正比厚度三次方，但可打印范围有限）；统一取向的各向异性把 D4 降到 D2，所以图案能偏离均匀板、却仍保留镜面对称。需要勘正的是撤销 $\theta$ 的原因：不是"自由度太多让程序混乱"，而是【代码核实】W10 工作在离共振区，惯性项压倒刚度项约 5 个数量级（$\omega^2 M \gg K$），$\theta$ 只通过刚度进入方程，\textbf{梯度近乎为零}---它是物理上推不动的寄生变量，随机漂移还污染下游。10 次对照实验证实：各向同性 PLA 在 4/6 个目标上赢过带"优化 $\theta$"的碳纤维板。便宜的 PLA 反而是推荐材料。

\subsection{评分之路："只在中间堆沙"的再归因}

最早的评分（沙近目标加分、远离扣分）总是产出"沙堆在中间"的设计，当时归因为"方程没有目标形状的解"。这只对了一半：结合木星机理，早期"中间堆沙"相当一部分是\textbf{中心伪影}---建模把中心标成最适合积沙的地方。物理限制是真的，但它的表现是"图案偏对称、偏简单"，不是"沙挤到中心"。前传也值得记录：三指标之前还有一代 IoU/Dice 逐像素评分，因离散不可导、模态一切换就跳变而失败；是"撒粉物理"的建模转换救活了评分系统---这也是方向 4 敢于在连续沙密度场上重新引入几何比对的依据。

\subsection{surrogate 到底是什么}

它不是现成工具，而是自写的板振动求解器（Kirchhoff--Love 方程离散到 $25\times 25$ 粗网格），用 PyTorch 实现的关键是\textbf{自动求导}---让"沙图对厚度的导数"可以直接算，梯度搜索才成立。要记住它的性格：\textbf{系统性乐观}（承诺 8--18 倍的提升，COMSOL 只兑现 1.5--4 倍），把板当薄纸、忽略三维效应。项目准则"用代理选方向，绝不用它宣布胜利"现在要再加一句：\textbf{COMSOL 也不能宣布胜利---只有真实的板可以。}

\subsection{离共振的决定与它的双重连锁反应}

"跳出特征模态"（不局限于共振频率）当时被金属板实验部分支持，幸运频率 165~Hz 是它的直接产物。现在看，这个决定有\textbf{两重}当时没人意识到的连锁反应：第一重，离共振区 $\omega^2 M \gg K$，让 $\theta$ 变成寄生变量（见 8.1 节）；第二重也是更致命的，\textbf{离共振工作点本身物理上不可复现}（响应弱、对模态偏移敏感），165~Hz 偏离最近本征模 2.87 个带宽---\textbf{"幸运频率"其实是病征，不是法宝}。值得注意的是，科普版第 7 节早就记下了相反的实验观察："要精确踩在那个音上---稍微偏离固有音高，中心推力就喧宾夺主毁掉图案。"实验和理论其实一直在说同一件事，只是当时没人把它和离共振路线的矛盾对起来。

\subsection{勘误的勘误：材料错配只是孤例}

上一轮分析曾把"优化器与 COMSOL 材料参数不一致"误判为第二个系统性根因。普查 32 对后更正：只有 1 对不一致，是一个过时的旧文件（基础 production\_design：优化器各向同性 vs COMSOL $sr=1.67$），其余全部一致。保留这条更正，正是"可信度分级"必要性的示范：把估计当定论，纠错成本会摊到后人头上。

\section{实施路线图}

\begin{table}[h]
\centering
\small
\begin{tabular}{@{}clll@{}}
\toprule
阶段 & 方向 & 工作量 & 预期收益 \\
\midrule
0 & 方向 0：实测标定回路（打印 + 敲击/扫频 + 拟合） & 中 & 裁判可信，所有数字落地 \\
1 & 方向 1：可复现性硬约束 + 本征频率离散选择 & 中 & 设计在真实板上震得出来 \\
1 & 方向 2a：可达目标库 + 裁决前置 & 小 & 不再向用户承诺不可能 \\
2 & 方向 2b：改物理（多激振器 / 边缘激励 / 非对称板） & 大 & 真正解锁非对称图案 \\
2 & 方向 3：统一激励模型 + surrogate 对照 COMSOL & 中 & 选板与验收是同一物理 \\
3 & 方向 4：评分重构 + 几何损失 + 随机重启 & 中 & 裁判可信后再磨尺子 \\
--- & 中心伪影 & 已结案 & mask 已上线；"正修"勿做 \\
\bottomrule
\end{tabular}
\caption{终版路线图。和初版的关键区别：先让仿真接触现实（阶段 0--1），再优化仿真内部的精度（阶段 3）。}
\end{table}

\section{结语}

本项目最大的方法论缺陷，是\textbf{全程在仿真里闭环：把 COMSOL 当真理，却从未让真理面对现实}。初版报告（包括它对评分与优化器的全部分析）都是在这个闭环\textbf{内部}做的改良---内容大多正确，位置全部靠后。留给接手者的三句话：\textbf{先建实验验证回路；把物理可复现性当硬约束；诚实承认单点激励的方形板做不出任意非对称图案}---然后，要么选物理允许的图案做到极致，要么动手改变物理。试错史的元教训也再次应验：最有价值的进展，从来都来自对旧结论的重新归因---包括对本报告自己的。

\end{document}
